\section{Approximate parent sensitivity and learned inhibitory routing}
\label{note:optional_extensions}

These two studies ask whether the nonlinear-population result requires an exact parent-response derivative and whether the context-to-inhibition map must be supplied. They use the same sixteen-neuron $[4,2]$ conductance population and interaction target as main Fig.~6, with the changes below. The two studies have separate development and confirmation seeds.

\paragraph{Design, selection and statistics.}
Each study used three development seeds to select an Adam rate independently for each rule from $0.01$, $0.03$ and $0.1$, minimizing mean validation NMSE. Twenty fresh paired seeds per study then compared the selected rates and the common rate $0.03$. Coincident rate policies share one trajectory; they are not independent replications. Training used 2,048 examples, minibatches of 128 and 4,096 updates, with validation on 1,024 examples every 128 updates, including initialization. The retained checkpoint minimizes validation NMSE; ordinary test uses 4,096 independently generated examples. Conductance log bounds $[-9,9]$, global gradient clipping at 10, initialization, masks, readout and loss normalization follow the original rescue experiment. Development did not evaluate test data. Separate generators keep minibatch sampling unchanged by proxy noise or shuffling.

The common-rate sensitivity contrasts compare the four-bin and noise-SD-$0.5$ rules with resistance gating and their respective shuffled controls. The selected-rate routing contrasts compare locally learned augmented routing with uniform routing and with locally learned resistance gating. These six comparisons were fixed before confirmation; two proposed recurrent contrasts were removed when that study was deferred without training. Paired whole-seed bootstrap intervals use 20,000 resamples. Exact two-sided sign-flip tests use one Holm correction across the six retained contrasts. Other comparisons and stress tests are descriptive. Full development outcomes, validation histories, source and checkpoint hashes, rate choices and seed-level outcomes are retained.

\paragraph{Noisy and coarse parent-state signals.}
Only the terminal learning signal changes; the forward network and target are unchanged. The noisy estimate substitutes $1-\tanh^2(V_p+\sigma\xi)$ for $f'_p(V_p)$, with independent standard Gaussian $\xi$ for every example, parent and update, and $\sigma\in\{0.25,0.5,1\}$ in normalized model-voltage units. It is not noise added to the forward voltage. The two-bin estimate uses the threshold $V_p=0.5$ and delivers the derivative at the midpoint of the occupied band. The four-bin estimate uses thresholds $0.25$, $0.5$ and $0.75$, with midpoint values $0.125$, $0.375$, $0.625$ and $0.875$. These bands span the physical range $0\leq V_p<1$ of this positive-conductance shunting model: terminal voltage is $E/(1+E+I)$, and parent voltage is a conductance-weighted sum of child voltages divided by leak, coupling and inhibitory conductances. The inhibitory attenuation $h$ is retained and the omitted coupling factor $\kappa$ is not supplied.

Shuffled controls permute the noisy or four-bin sensitivity independently within each parent and context group of a minibatch, preserving the empirical sensitivity distribution while disrupting its assignment to examples. A context-mean control replaces each four-bin sensitivity by its parent/context minibatch mean. The latter retains the first moment but removes within-group variation. Random permutations can retain fixed points, and minibatch means include the current example; neither finite-minibatch control is asserted to remove every trace of sample dependence in expectation. Both are distinct from the cue-only resistance rule.

An initial development pilot used bin thresholds extending beyond the reachable voltage range. This made its nominal two-bin estimate constant and left upper bins unused. The pilot was stopped, all 84 completed fits and executable sources were retained, and all development fits were repeated with the bins corrected to the physical range before any confirmation test was evaluated. The amendment is part of the archived protocol. The pilot is not used as evidence against coarse sensitivities that vary with the input.

Both voltage resolutions and all three noise levels retained low ordinary-test error, close to the unperturbed augmented rule (main Fig.~6E; Fig.~\ref{fig:si_optional_extensions}A). Shuffling the estimates across examples, or replacing the four-bin estimate by its context mean, left a large error. Thus, in this task the useful example dependence did not require an exact numerical derivative. The primary contrasts and uncertainty estimates are reported in Table~\ref{tab:optional_extensions}; the full selected-rate comparisons and distractor-stress outcomes remain in Source Data.

\paragraph{Learning the cue-to-inhibition map.}
A shared $4\times4$ matrix $W$ maps the supplied one-hot context cue $\mathbf e_c$ to relative disinhibition,
\[
\mathbf q(c)=\operatorname{softmax}(\mathbf e_c^{\mathsf T}W),
\qquad a_b^{\rm I}(c)=4\{1-q_b(c)\}.
\]
The total inhibitory activity is therefore 12 in every context; total inhibitory conductance can still change as the original synaptic weights learn. Each matrix entry starts as an independent zero-mean Gaussian with standard deviation $0.02$. The matrix is optimized with the same Adam rate as the conductance network and clamped to $[-9,9]$. Every routing condition sets the terminal excitatory cue inputs to zero, so the cue reaches the network only through inhibition. Stream-specific sensory contacts remain fixed.

The local controller update sums parent-local parameter eligibilities weighted by each neuron's supplied somatic error. It does not receive an explicit desired route or the exact downstream coupling factors. The augmented rule supplies $h f'_p$ to terminal synapses; the resistance comparison supplies $h$. Because a controller weight affects several parents, its local update need not be a constant multiple of the exact gradient. A separate reference gives only the controller the exact task gradient, keeping augmented local credit in the conductance network. Supplied routing uses $\mathbf q=\mathbf e_c$, uniform routing uses $q_b=1/4$, and wrong routing cyclically shifts the supplied route. These controls share the same removal of terminal cue contacts. Uniform routing consequently removes context information from the forward computation; it is a forward-selection control, not an information-matched learning-rule comparison. The controller is a small normalized map receiving already identified contexts, not a model of a distributed inhibitory circuit or autonomous context discovery.

The controller learned selective routing with either terminal rule, but this did not remove the need for input-dependent terminal credit. Augmented learning reached low error, while resistance gating retained a substantial interaction deficit despite selecting the matching stream (main Fig.~6F; Fig.~\ref{fig:si_optional_extensions}C,D). The exact-controller reference also reached low error. This distinguishes acquisition of branch selection from supplying the additional parent-state dependence used by terminal updates. Selected- and common-rate outcomes, all controller maps and selected-state bound occupancy are included in Source Data. Controller bound occupancy describes selected checkpoints; the training-contact count in the endpoint table concerns core conductances, not transient controller-logit contacts.

\paragraph{Scope.}
These studies relax precision of the parent-state estimate and prior specification of a cue-to-route map in the tested model. They retain Adam, fixed sensory grouping and supplied neuron-specific somatic errors. Approximate state dependence does not establish a biological carrier for the parent-to-terminal signal. The source records are under \path{source_data/optional_extensions/}; executable mechanisms and frozen protocols are documented in \path{code/optional_extensions/}.

\clearpage
\begin{figure}[p]
\centering
\includegraphics[width=\textwidth]{supplementary/curated/optional_extensions.pdf}
\caption{\textbf{Noise, target and routing controls delimit the nonlinear-population result.}
\textbf{A}, Parent-sensitivity cohort: ordinary-test error when noise is added to the voltage used to estimate sensitivity, with no change to forward voltage. Zero noise is the augmented rule. Horizontal lines show mean resistance-gate error and the shuffled noise-SD-$0.5$ control; these means nearly coincide. Adam rate is 0.03. Main Fig.~6E shows the coarse estimates and matched shuffled controls.
\textbf{B}, Original rescue cohort: interaction target (circles) versus no interaction (squares), retaining the same nonlinear parents and common Adam rate 0.03. Exact, broadcast, resistance, augmented and within-parent/context derivative-shuffle rules share twenty paired seeds. Removing the interaction removes the large resistance-gate deficit. These controls are specified in Section~S5; selected-rate interaction outcomes appear in main Fig.~6C.
\textbf{C}, Separate routing cohort at common rate 0.03, with direct terminal cue contacts removed in every condition. Learned local and exact-gradient controllers both accompany augmented terminal credit; the resistance condition changes that terminal signal. Supplied, uniform and cyclically wrong routing are forward controls. Main Fig.~6F shows the selected-rate comparison.
\textbf{D}, Mean route probabilities of locally trained controllers at selected rates, with augmented (left) or resistance (right) terminal credit. Rows identify the cue and columns the disinhibited stream; probabilities sum to one per row. These maps describe routing, not a second performance endpoint.
Dots show all twenty seeds per cohort; large markers and intervals show means and 95\% whole-seed bootstrap intervals. The three cohorts are separate; rate policies reuse seeds. Source Data retain every condition, including exact-credit and context-mean proxy controls, distractor stress, validation histories and all prespecified contrasts.}
\label{fig:si_optional_extensions}
\end{figure}
\clearpage
